\section{A simple model}

Treat the sheet as a thin, uniform medium and let $c(r,t)$
be the concentration of pigment at distance $r$ from where
the drop landed. Within the wet region the pigment obeys
the diffusion equation
\begin{equation}
  \frac{\partial c}{\partial t} = D\,\nabla^2 c,
  \label{eq:diffusion}
\end{equation}
where $D$ is the diffusion coefficient of the pigment in
the paper. For a drop that deposits a mass $M$ at a single
point, the solution of~\eqref{eq:diffusion} in two
dimensions is
\begin{equation}
  c(r,t) = \frac{M}{4\pi D t}\,
           \exp\!\left(-\frac{r^2}{4 D t}\right).
  \label{eq:solution}
\end{equation}
The profile keeps its shape and only widens: its width is
set by the length $\ell(t) = \sqrt{4 D t}$.

The water front moves in the same way. For capillary flow
through channels of radius $a$, the Washburn equation
gives the distance reached after time $t$ as
\begin{equation}
  L(t) = \sqrt{\frac{\gamma\, a\, t \cos\theta}{2\eta}},
  \label{eq:washburn}
\end{equation}
with surface tension $\gamma$, contact angle $\theta$, and
viscosity $\eta$. Both lengths grow as $\sqrt{t}$, which
is why a stain spreads quickly at first and then seems to
stop.
